% ============================================================
%  GAN Presentation — Maciej Leśniak & Bartłomiej Obrochta
%  Sztuczne Sieci Neuronowe | 22 kwietnia 2026
%  Compile: xelatex gan_presentation.tex (twice)
% ============================================================
\documentclass[aspectratio=169,12pt]{beamer}

\usepackage{fontspec}
\usepackage{polyglossia}
\setmainlanguage{polish}
\setmainfont{DejaVu Serif}
\setsansfont{DejaVu Sans}
\setmonofont{Liberation Mono}

\usepackage{amsmath,amssymb,mathtools}
\usepackage{tikz}
\usepackage{pgfplots}
\pgfplotsset{compat=1.18}
\usepackage{pgf-pie}
\usepackage{xcolor}
\usepackage{booktabs}
\usepackage{graphicx}
\usetikzlibrary{arrows.meta,positioning,shapes.geometric,fit,calc,backgrounds}

% ── Colour palette ───────────────────────────────────────────
\definecolor{ganblue}{RGB}{15,52,96}
\definecolor{ganaccent}{RGB}{232,93,16}
\definecolor{ganteal}{RGB}{0,168,150}
\definecolor{ganlightbg}{RGB}{240,245,255}
\definecolor{ganmid}{RGB}{100,140,200}
\definecolor{gangray}{RGB}{80,90,110}

% ── Theme ────────────────────────────────────────────────────
\usetheme{metropolis}
\setbeamercolor{palette primary}{bg=ganblue,fg=white}
\setbeamercolor{titlelike}{parent=palette primary}
\setbeamercolor{frametitle}{bg=ganblue,fg=white}
\setbeamercolor{progress bar}{fg=ganaccent}
\setbeamercolor{alerted text}{fg=ganaccent}
\setbeamercolor{block title}{bg=ganblue!90,fg=white}
\setbeamercolor{block body}{bg=ganlightbg}
\setbeamercolor{structure}{fg=ganblue}
\setbeamerfont{frametitle}{size=\large,series=\bfseries}
\setbeamerfont{title}{size=\huge,series=\bfseries}
\metroset{progressbar=frametitle,sectionpage=none,numbering=fraction}

% ── Helpers ──────────────────────────────────────────────────
\newcommand{\imgph}[2][.92\linewidth]{%
  \begin{tikzpicture}
    \node[draw=ganmid,dashed,fill=ganlightbg,minimum width=#1,
          minimum height=3.6cm,rounded corners=6pt,inner sep=0](b){};
    \node[text=ganmid,font=\small\itshape] at (b.center){[ #2 ]};
  \end{tikzpicture}}

\newcommand{\abox}[1]{%
  \begin{tikzpicture}
    \node[draw=ganaccent,fill=ganaccent!8,rounded corners=4pt,
          inner sep=7pt,text width=\linewidth-18pt,line width=1.4pt,font=\small]{#1};
  \end{tikzpicture}}

% ── Metadata ─────────────────────────────────────────────────
\title{Generative Adversarial Networks}
\subtitle{Sieci neuronowe typu GAN}
\author{Maciej Leśniak \and Bartłomiej Obrochta}
\institute{Sztuczne Sieci Neuronowe — AGH Kraków}
\date{22 kwietnia 2026}

\begin{document}

{\setbeamertemplate{footline}{}
\begin{frame}[plain]\maketitle\end{frame}}

%% ─────────────────────── AGENDA ─────────────────────────────
\begin{frame}{Plan prezentacji}
\begin{columns}[T]
\begin{column}{0.48\textwidth}
\begin{enumerate}
  \item \textbf{Rys historyczny}
  \item Intuicja: fałszerz i policjant
  \item Architektura GAN
  \item \textbf{Formalizm matematyczny}
  \item Trening krok po kroku
  \item \textbf{Problemy} (mode collapse, gradienty)
\end{enumerate}
\end{column}
\begin{column}{0.48\textwidth}
\begin{enumerate}\setcounter{enumi}{6}
  \item Warianty: DCGAN, CGAN, WGAN, StyleGAN, CycleGAN
  \item Zastosowania
  \item GAN vs.\ Diffusion Models
  \item Etyka i deepfake'i
  \item Live demo — Jupyter
\end{enumerate}
\end{column}
\end{columns}
\vspace{0.6em}
\abox{Pytania możecie zadawać na bieżąco.}
\end{frame}

%% ─────────────────────── HISTORIA ───────────────────────────
\begin{frame}{Rys historyczny}
\begin{columns}[T]
\begin{column}{0.52\textwidth}
\begin{tikzpicture}[
  box/.style={draw=ganblue!40,fill=ganlightbg,rounded corners=4pt,
    minimum width=4.3cm,minimum height=0.62cm,font=\small,inner sep=4pt},
  acc/.style={box,fill=ganaccent!15,draw=ganaccent},
  yr/.style={font=\bfseries\small,text=ganaccent},
  a/.style={ganblue!50,line width=1.1pt}]
  \node[yr](y0){2000s};
  \node[box,right=0.3cm of y0](b0){RBM + MCMC (Hinton)};
  \node[yr,below=0.32cm of y0](y1){2013};
  \node[box,right=0.3cm of y1](b1){VAE — rozmyte obrazy};
  \node[yr,below=0.32cm of y1](y2){2014};
  \node[acc,right=0.3cm of y2](b2){\textbf{GAN — Goodfellow}};
  \node[yr,below=0.32cm of y2](y3){2015};
  \node[box,right=0.3cm of y3](b3){DCGAN};
  \node[yr,below=0.32cm of y3](y4){2017};
  \node[box,right=0.3cm of y4](b4){WGAN, Pix2Pix, CycleGAN};
  \node[yr,below=0.32cm of y4](y5){2018};
  \node[box,right=0.3cm of y5](b5){StyleGAN (NVIDIA)};
  \node[yr,below=0.32cm of y5](y6){2020+};
  \node[box,right=0.3cm of y6](b6){Diffusion Models dominują};
  \foreach \x/\y in {y0/y1,y1/y2,y2/y3,y3/y4,y4/y5,y5/y6}
    \draw[a](\x.south)--(\y.north);
\end{tikzpicture}
\end{column}
\begin{column}{0.45\textwidth}
\begin{block}{Ian Goodfellow, 2014}
\small Prototyp zakodowany \textbf{w jedną noc}. Zadziałało od razu.
\end{block}
\vspace{0.4em}
\abox{\textbf{Yann LeCun (2016):}\\
\emph{„GANy to najciekawszy pomysł w~ML ostatnich dziesięciu lat."}}
\vspace{0.5em}
\imgph[.95\linewidth]{Zdjęcie: Ian Goodfellow}
\end{column}
\end{columns}
\end{frame}

%% ─────────────────────── INTUICJA ───────────────────────────
\begin{frame}{Intuicja: fałszerz i policjant}
\begin{columns}[T]
\begin{column}{0.55\textwidth}
\begin{tikzpicture}[
  bl/.style={rounded corners=6pt,minimum width=2.5cm,minimum height=0.95cm,
    align=center,font=\small\bfseries},
  a/.style={-{Stealth},thick,line width=1.2pt}]
  \node[bl,fill=ganaccent!80,text=white](gen){Generator\\(fałszerz)};
  \node[bl,fill=ganblue,text=white,right=2.5cm of gen](dis){Discriminator\\(policjant)};
  \node[bl,fill=ganteal!70,text=white,above=1.0cm of dis](real){Prawdziwe dane};
  \node[font=\scriptsize,text=gangray,below=0.5cm of gen](noise){szum $\mathbf{z}\sim\mathcal{N}(0,\mathbf{I})$};
  \node[font=\scriptsize\bfseries,text=ganblue,right=0.5cm of dis](out){Real / Fake};
  \draw[a,gangray!60](noise)--(gen);
  \draw[a,ganaccent!90](gen)--node[above,font=\scriptsize]{$G(\mathbf{z})$}(dis);
  \draw[a,ganteal!90](real)--(dis);
  \draw[a,ganblue](dis)--(out);
  \draw[a,dashed,gangray,bend right=40](dis.south) to
    node[below,font=\scriptsize,text=gangray]{gradient / feedback}(gen.south);
\end{tikzpicture}
\end{column}
\begin{column}{0.42\textwidth}
\begin{itemize}
  \item Generator \textbf{nigdy nie widzi} prawdziwych danych bezpośrednio
  \item Jedyna informacja = ocena Discriminatora
  \item To \textbf{czysta rywalizacja} — nie współpraca
  \item Rywalizacja sprawia, że oba modele się doskonalą
\end{itemize}
\vspace{0.4em}
\imgph[.95\linewidth]{Ilustracja: fałszerz i policjant}
\end{column}
\end{columns}
\end{frame}

%% ─────────────────────── ARCHITEKTURA ───────────────────────
\begin{frame}{Architektura GAN}
\begin{center}
\begin{tikzpicture}[
  bl/.style={rounded corners=5pt,minimum width=2.1cm,minimum height=0.88cm,
    align=center,font=\small},
  a/.style={-{Stealth[scale=1.1]},line width=1.2pt}]
  \node[bl,fill=gangray!20](z){$\mathbf{z}$\\\tiny szum};
  \node[bl,fill=ganaccent!80,text=white,right=1.2cm of z](G){\textbf{Generator}\\\tiny $G_\theta(\mathbf{z})$};
  \node[bl,fill=ganblue!70,text=white,right=1.5cm of G](D){\textbf{Discriminator}\\\tiny $D_\phi(\cdot)$};
  \node[bl,fill=ganteal!60,text=white,above=0.95cm of D](X){Prawdziwe\\\tiny $\mathbf{x}\sim p_{\text{data}}$};
  \node[font=\scriptsize\bfseries,text=ganblue,right=0.7cm of D](out){$D\in[0,1]$\\Real/Fake};
  \draw[a,gangray!60](z)--(G);
  \draw[a,ganaccent!90](G)--node[above,font=\tiny]{$G(\mathbf{z})$}(D);
  \draw[a,ganteal!90](X)--(D);
  \draw[a,ganblue](D)--(out);
  \draw[a,dashed,ganaccent!70,bend right=32](D.south) to
    node[below,font=\tiny,text=gangray]{$\nabla_\theta\mathcal{L}_G$}(G.south);
  \node[font=\tiny,text=red!70,above=0.15cm of G.north east]{G nie widzi $\mathbf{x}$!};
\end{tikzpicture}
\end{center}
\vspace{0.2em}
\begin{columns}[T]
\begin{column}{0.48\textwidth}
\textbf{Generator} $G_\theta$:
\begin{itemize}
  \item Wejście: $\mathbf{z}\in\mathbb{R}^{100}$, $\mathbf{z}\sim\mathcal{N}(0,\mathbf{I})$
  \item Wyjście: obraz ($28{\times}28$ — MNIST)
  \item Na początku: wizualny szum
\end{itemize}
\end{column}
\begin{column}{0.48\textwidth}
\textbf{Discriminator} $D_\phi$:
\begin{itemize}
  \item Klasyfikator binarny
  \item Wyjście: $D(\mathbf{x})\in[0,1]$
  \item $1=$ prawdziwy,\ $0=$ fałszywy
\end{itemize}
\end{column}
\end{columns}
\end{frame}

%% ─────────────────────── MATH ───────────────────────────────
\begin{frame}{Formalizm matematyczny — Minimax}
\begin{center}
\begin{tikzpicture}
\node[draw=ganblue,fill=ganlightbg,rounded corners=6pt,inner sep=12pt,line width=1.8pt]{
$\displaystyle
\min_{G}\;\max_{D}\;V(D,G)=
\underbrace{\mathbb{E}_{\mathbf{x}\sim p_{\text{data}}}\!\bigl[\log D(\mathbf{x})\bigr]}_{\text{prawdziwe dane}}
+
\underbrace{\mathbb{E}_{\mathbf{z}\sim p_z}\!\bigl[\log\!\bigl(1-D(G(\mathbf{z}))\bigr)\bigr]}_{\text{wygenerowane dane}}
$};
\end{tikzpicture}
\end{center}
\begin{columns}[T]
\begin{column}{0.48\textwidth}
\textbf{Discriminator} ($\max_D$):
\begin{itemize}
  \item Chce $D(\mathbf{x})\to 1$ (prawdziwe)
  \item Chce $D(G(\mathbf{z}))\to 0$ (fałszywe)
  \item Maksymalizuje $V$
\end{itemize}
\textbf{Generator} ($\min_G$):
\begin{itemize}
  \item Chce $D(G(\mathbf{z}))\to 1$
  \item Minimalizuje $V$
\end{itemize}
\textbf{Nash optimum:} $D^*(\mathbf{x})=\tfrac{1}{2}\ \forall\,\mathbf{x}$
\end{column}
\begin{column}{0.49\textwidth}
\textbf{Non-saturating loss (trik):}
\vspace{0.3em}

Oryginalna (nasyca się):
\[\min_G\log\!\bigl(1-D(G(\mathbf{z}))\bigr)\]
Zamiennik (silniejszy gradient):
\begin{center}
\begin{tikzpicture}
\node[draw=ganaccent,fill=ganaccent!10,rounded corners=4pt,inner sep=8pt]{
$\displaystyle\max_G\;\log D(G(\mathbf{z}))$};
\end{tikzpicture}
\end{center}
\vspace{0.2em}
{\small Matematycznie ekwiwalentne w optimum, ale gradient na początku treningu dużo silniejszy.}
\end{column}
\end{columns}
\end{frame}

%% ─────────────────────── LOSS PLOTS ─────────────────────────
\begin{frame}{Funkcje straty — wizualizacja}
\begin{columns}[T]
\begin{column}{0.48\textwidth}
\begin{tikzpicture}
\begin{axis}[width=6.4cm,height=5.5cm,
  xlabel={$D(G(\mathbf{z}))$},ylabel={wartość},
  xmin=0.01,xmax=0.99,ymin=-4,ymax=0.4,
  legend style={font=\tiny,at={(0.97,0.97)},anchor=north east},
  title={\small Strata Generatora},title style={font=\small\bfseries},
  grid=major,grid style={gray!20}]
  \addplot[ganaccent,thick,domain=0.01:0.99,samples=120]{ln(1-x)};
  \addlegendentry{$\log(1-D(G))$ [nasyca się]}
  \addplot[ganteal,thick,dashed,domain=0.01:0.99,samples=120]{ln(x)};
  \addlegendentry{$\log D(G)$ [NS, lepszy]}
\end{axis}
\end{tikzpicture}
\end{column}
\begin{column}{0.48\textwidth}
\begin{tikzpicture}
\begin{axis}[width=6.4cm,height=5.5cm,
  xlabel={iteracja},ylabel={loss},
  xmin=0,xmax=100,ymin=0,ymax=3.2,
  legend style={font=\tiny,at={(0.97,0.97)},anchor=north east},
  title={\small Typowy przebieg treningu GAN},title style={font=\small\bfseries},
  grid=major,grid style={gray!20}]
  \addplot[ganblue,thick,domain=0:100,samples=150]
    {1.2+0.5*sin(deg(x/4.5))*exp(-x/50)+0.08*sin(deg(x*1.7))};
  \addlegendentry{$\mathcal{L}_D$}
  \addplot[ganaccent,thick,domain=0:100,samples=150]
    {1.8-0.4*sin(deg(x/3+1))*exp(-x/40)+0.1*sin(deg(x*2.1))};
  \addlegendentry{$\mathcal{L}_G$}
  \addplot[ganteal,dashed,thick,domain=0:100,samples=2]{2.772};
  \addlegendentry{Nash: $2\ln 2$}
\end{axis}
\end{tikzpicture}
\end{column}
\end{columns}
\end{frame}

%% ─────────────────────── TRAINING ALGO ──────────────────────
\begin{frame}{Algorytm treningu GAN — krok po kroku}
\begin{columns}[T]
\begin{column}{0.50\textwidth}
\begin{block}{Krok 1 — trenuj Discriminatora}
\begin{enumerate}\small
  \item Zamroź $G$
  \item Real batch $\{\mathbf{x}^{(i)}\}$, etykieta $y=1$
  \item Fake batch $\{G(\mathbf{z}^{(i)})\}$, etykieta $y=0$
  \item $\mathcal{L}_D = {-}\tfrac{1}{m}\!\sum_i\bigl[\log D(\mathbf{x}^{(i)})
    {+}\log(1{-}D(G(\mathbf{z}^{(i)})))\bigr]$
  \item \texttt{.detach()} na $G(\mathbf{z})$ \alert{— kluczowe!}
  \item Backprop, aktualizuj tylko $\phi$
\end{enumerate}
\end{block}
\end{column}
\begin{column}{0.48\textwidth}
\begin{block}{Krok 2 — trenuj Generator}
\begin{enumerate}\small
  \item Zamroź $D$
  \item Fake batch \emph{bez} \texttt{.detach()}
  \item Etykieta $y=\alert{1}$ (NIE 0!)
  \item $\mathcal{L}_G = {-}\tfrac{1}{m}\!\sum_i \log D(G(\mathbf{z}^{(i)}))$
  \item Backprop przez $D$ do $G$
  \item Aktualizuj tylko $\theta$
\end{enumerate}
\end{block}
\vspace{0.4em}
\abox{\textbf{Osobne optymizatory:} Adam, $\eta{=}0.0002$, $\beta_1{=}0.5$}
\end{column}
\end{columns}
\end{frame}

%% ─────────────────────── MODE COLLAPSE ──────────────────────
\begin{frame}{Problem: Mode Collapse}
\begin{columns}[T]
\begin{column}{0.50\textwidth}
\begin{tikzpicture}
\begin{axis}[width=6.4cm,height=5cm,
  xlabel={$x$},ylabel={$p(x)$},
  xmin=-6,xmax=6,ymin=0,ymax=0.48,
  title={\small Rozkład prawdziwy vs.\ collapsed},title style={font=\small\bfseries},
  legend style={font=\tiny,at={(0.97,0.97)},anchor=north east},
  grid=major,grid style={gray!15},axis lines=left]
  \addplot[ganteal,thick,domain=-6:6,samples=200]
    {0.21*exp(-0.5*(x+3)^2/0.55)+0.21*exp(-0.5*(x-3)^2/0.55)};
  \addlegendentry{$p_{\text{data}}$ — bimodalny}
  \addplot[ganaccent,thick,dashed,domain=-6:6,samples=200]
    {0.42*exp(-0.5*(x-3)^2/0.25)};
  \addlegendentry{$p_G$ — collapsed!}
\end{axis}
\end{tikzpicture}
\end{column}
\begin{column}{0.47\textwidth}
\begin{itemize}
  \item Generator „kolapsuje" do 1–2 trybów zamiast modelować cały rozkład
  \item Przykład: MNIST — same jedynki
\end{itemize}
\vspace{0.5em}
\textbf{Metody zaradcze:}
\begin{itemize}
  \item \textbf{Minibatch discrimination} — Discriminator porównuje próbki w batchu
  \item \textbf{Feature matching} — dopasowanie statystyk pośrednich warstw
  \item \textbf{WGAN} — fundamentalna zmiana metryki
\end{itemize}
\vspace{0.4em}
\imgph[.95\linewidth]{Mode collapse — przykład na MNIST}
\end{column}
\end{columns}
\end{frame}

%% ─────────────────────── GRADIENTY + FID ────────────────────
\begin{frame}{Zanikające gradienty i metryki jakości}
\begin{columns}[T]
\begin{column}{0.46\textwidth}
\textbf{Zanikające gradienty:}
\begin{tikzpicture}
\begin{axis}[width=6cm,height=4.2cm,
  xlabel={$D(G(\mathbf{z}))$},ylabel={gradient},
  xmin=0.001,xmax=0.5,ymin=-15,ymax=0,
  title={\small $\nabla\log(1-D(G(\mathbf{z})))$},
  title style={font=\scriptsize\bfseries},
  grid=major,grid style={gray!15},axis lines=left]
  \addplot[ganaccent,thick,domain=0.001:0.49,samples=150]{-1/(1-x)};
\end{axis}
\end{tikzpicture}
{\small Zbyt dobry $D$ $\Rightarrow$ gradient $\approx 0$ $\Rightarrow$ G~przestaje się uczyć}
\end{column}
\begin{column}{0.51\textwidth}
\textbf{Metryki jakości:}
\vspace{0.4em}

\begin{tabular}{@{}lll@{}}
\toprule Metryka & Rok & Idea \\ \midrule
IS & 2016 & Inception Score \\
\textbf{FID} & \textbf{2017} & \textbf{Fréchet Inception Dist.} \\
\bottomrule
\end{tabular}

\vspace{0.5em}
\begin{block}{FID (aktualny standard)}
\[
\mathrm{FID} = \|\mu_r{-}\mu_g\|^2
+ \mathrm{tr}\!\Bigl(\Sigma_r{+}\Sigma_g
  {-}2\bigl(\Sigma_r\Sigma_g\bigr)^{\!\frac{1}{2}}\Bigr)
\]
\textbf{Mniejsze FID = lepszy model.}
\end{block}
\end{column}
\end{columns}
\end{frame}

%% ─────────────────────── DCGAN ──────────────────────────────
\begin{frame}{DCGAN (2015) — Deep Convolutional GAN}
\begin{columns}[T]
\begin{column}{0.50\textwidth}
\textbf{Kluczowe innowacje:}
\begin{itemize}
  \item \textbf{Transposed convolutions} w G:\\
        $4{\times}4\to8{\times}8\to\cdots\to64{\times}64$
  \item \textbf{Strided convolutions} zamiast MaxPool w D
  \item Batch Normalization po każdej warstwie
  \item ReLU (G) / LeakyReLU (D) / Tanh (wyjście G)
\end{itemize}
\vspace{0.4em}
\textbf{Arytmetyka w przestrzeni latentnej:}
\vspace{0.3em}

\begin{tikzpicture}[font=\small,node distance=0.15cm]
  \node(a){$\mathbf{z}_{\text{mężczyzna+okulary}}$};
  \node[below=0.1cm of a](b){$-\;\mathbf{z}_{\text{mężczyzna}}$};
  \node[below=0.1cm of b](c){$+\;\mathbf{z}_{\text{kobieta}}$};
  \draw[ganblue,thick]([xshift=-3pt]a.south west)--([xshift=-3pt]c.north west);
  \node[right=0.3cm of c,text=ganaccent,font=\bfseries\small]{$=$ kobieta+okulary};
\end{tikzpicture}
\end{column}
\begin{column}{0.47\textwidth}
\imgph[.95\linewidth]{Architektura generatora DCGAN}
\vspace{0.5em}
\imgph[.95\linewidth]{Wygenerowane twarze DCGAN}
\end{column}
\end{columns}
\end{frame}

%% ─────────────────────── WGAN ───────────────────────────────
\begin{frame}{WGAN — Earth Mover's Distance}
\begin{columns}[T]
\begin{column}{0.50\textwidth}
\textbf{Problem z JS-dywergencją:}
\[
\mathrm{JSD}(p_r\|p_g)=
\begin{cases}\log 2 & \text{gdy }p_r\cap p_g=\emptyset\\
<\log 2 & \text{inaczej}\end{cases}
\]
Stała wartość $\Rightarrow$ brak gradientu.

\vspace{0.4em}
\textbf{Earth Mover's Distance:}
\[W(p_r,p_g)=\inf_{\gamma\in\Pi(p_r,p_g)}\mathbb{E}_{(x,y)\sim\gamma}\bigl[\|x{-}y\|\bigr]\]

\textbf{Obiektyw WGAN (Critic $f$):}
\[\min_G\max_{\|f\|_L\le1}\mathbb{E}_{x\sim p_r}[f(x)]-\mathbb{E}_{z}[f(G(z))]\]
\end{column}
\begin{column}{0.47\textwidth}
\begin{tikzpicture}
\begin{axis}[width=6.2cm,height=4.5cm,
  xlabel={odległość między rozkładami},ylabel={wartość metryki},
  xmin=0,xmax=5,ymin=0,ymax=4,
  legend style={font=\tiny,at={(0.06,0.95)},anchor=north west},
  title={\small JS vs.\ Wasserstein},title style={font=\small\bfseries},
  grid=major,grid style={gray!15}]
  \addplot[gangray,thick,domain=0:5,samples=3]{2.772};
  \addlegendentry{JSD = stała}
  \addplot[ganaccent,thick,domain=0:5,samples=50]{0.7*x};
  \addlegendentry{$W$ — liniowo rośnie}
\end{axis}
\end{tikzpicture}
\vspace{0.3em}
\abox{\textbf{Efekty WGAN:} stabilniejszy trening, brak mode collapse, loss Critica \textbf{koreluje z jakością}}
\end{column}
\end{columns}
\end{frame}

%% ─────────────────────── STYLEGAN ───────────────────────────
\begin{frame}{StyleGAN (NVIDIA, 2018)}
\begin{columns}[T]
\begin{column}{0.50\textwidth}
\begin{tikzpicture}[
  bl/.style={rounded corners=4pt,minimum width=3cm,minimum height=0.68cm,
    font=\small,align=center},
  a/.style={-{Stealth},thick}]
  \node[bl,fill=gangray!20](z){$\mathbf{z}\in\mathcal{Z}$};
  \node[bl,fill=ganblue!70,text=white,below=0.42cm of z](f){Mapping Net $f$ ($8{\times}$FC)};
  \node[bl,fill=ganteal!60,text=white,below=0.42cm of f](w){$\mathbf{w}\in\mathcal{W}$ (rozplątana)};
  \node[bl,fill=ganaccent!70,text=white,below=0.42cm of w](ada){AdaIN (każda warstwa konw.)};
  \node[bl,fill=ganlightbg,draw=ganblue,below=0.42cm of ada](img){Obraz wyjściowy};
  \draw[a](z)--(f); \draw[a](f)--(w); \draw[a](w)--(ada); \draw[a](ada)--(img);
\end{tikzpicture}
\vspace{0.3em}
{\small $\mathcal{Z}$: splątana $\;\to\;$ $\mathcal{W}$: niezależne atrybuty}
\end{column}
\begin{column}{0.47\textwidth}
\textbf{Kontrola stylu per warstwa:}
\begin{itemize}
  \item Niskie warstwy $\to$ kształt, poza
  \item Średnie warstwy $\to$ rysy twarzy
  \item Wysokie warstwy $\to$ kolor, tekstura
\end{itemize}
\vspace{0.4em}
\abox{\textbf{thispersondoesnotexist.com}\\Żadna z tych twarzy nie istnieje.\\StyleGAN działa w czasie rzeczywistym.}
\vspace{0.4em}
\imgph[.95\linewidth]{Twarze wygenerowane przez StyleGAN}
\end{column}
\end{columns}
\end{frame}

%% ─────────────────────── CYCLEGAN ───────────────────────────
\begin{frame}{CycleGAN (2017) — bez sparowanych danych}
\begin{columns}[T]
\begin{column}{0.50\textwidth}
\begin{center}
\begin{tikzpicture}[
  bl/.style={rounded corners=5pt,minimum width=1.8cm,minimum height=0.78cm,
    font=\small,align=center},
  a/.style={-{Stealth},thick,line width=1.2pt}]
  \node[bl,fill=ganteal!50,text=white](X){$X$\\koń};
  \node[bl,fill=ganaccent!70,text=white,right=2.5cm of X](Y){$Y$\\zebra};
  \node[bl,fill=ganblue!50,text=white,below=0.9cm of X](Xr){$F(G(x))\approx x$};
  \node[bl,fill=ganblue!50,text=white,below=0.9cm of Y](Yr){$G(F(y))\approx y$};
  \draw[a,ganteal](X) to[bend left=22] node[above,font=\tiny]{$G$}(Y);
  \draw[a,ganaccent](Y) to[bend left=22] node[below,font=\tiny]{$F$}(X);
  \draw[a,dashed,ganblue!60](Y)--(Yr);
  \draw[a,dashed,ganblue!60](X)--(Xr);
\end{tikzpicture}
\end{center}
\vspace{0.3em}
\textbf{Cycle Consistency Loss:}
\[
\mathcal{L}_{\mathrm{cyc}}=
\mathbb{E}_x\!\bigl[\|F(G(x)){-}x\|_1\bigr]
+\mathbb{E}_y\!\bigl[\|G(F(y)){-}y\|_1\bigr]
\]
\end{column}
\begin{column}{0.47\textwidth}
\textbf{Przykłady zastosowań:}
\begin{itemize}
  \item Konie $\leftrightarrow$ Zebry
  \item Lato $\leftrightarrow$ Zima
  \item Zdjęcia $\leftrightarrow$ Obrazy (Monet)
  \item Mapa $\leftrightarrow$ Zdjęcie satelitarne
\end{itemize}
\vspace{0.4em}
\imgph[.95\linewidth]{CycleGAN: koń $\to$ zebra}
\vspace{0.3em}
{\small\textbf{Przewaga:} dwa \emph{niesparowane} zbiory wystarczą}
\end{column}
\end{columns}
\end{frame}

%% ─────────────────────── ZASTOSOWANIA ───────────────────────
\begin{frame}{Zastosowania GANów}
\begin{columns}[T]
\begin{column}{0.32\textwidth}
\begin{block}{Przemysł}
\small\begin{itemize}
  \item Super-resolution (ESRGAN)
  \item Zdjęcia produktów
  \item Augmentacja danych
\end{itemize}
\end{block}
\end{column}
\begin{column}{0.32\textwidth}
\begin{block}{Medycyna}
\small\begin{itemize}
  \item Syntetyczne MRI/CT
  \item Prywatność pacjentów
  \item Rzadkie choroby
\end{itemize}
\end{block}
\end{column}
\begin{column}{0.32\textwidth}
\begin{block}{Nauka}
\small\begin{itemize}
  \item CERN — CaloGAN
  \item Drug discovery (MolGAN)
  \item Fizyka cząstek
\end{itemize}
\end{block}
\end{column}
\end{columns}
\vspace{0.5em}
\begin{center}
\begin{tikzpicture}
\begin{axis}[width=11cm,height=3.7cm,
  symbolic x coords={Super-res,Inpainting,Medycyna,{Gry RT},Drug disc.,{Deepfake det.}},
  xtick=data,xticklabel style={font=\small},
  ylabel={dojrzałość},ymin=0,ymax=10,
  ybar,bar width=0.9cm,
  title={\small Dojrzałość zastosowań GANów},title style={font=\small\bfseries},
  grid=major,grid style={gray!15},
  nodes near coords,nodes near coords style={font=\tiny,above}]
  \addplot[fill=ganblue!80,draw=none]
    coordinates{(Super-res,9)(Inpainting,8)(Medycyna,6)({Gry RT},8)(Drug disc.,5)({Deepfake det.},7)};
\end{axis}
\end{tikzpicture}
\end{center}
\end{frame}

%% ─────────────────────── GAN vs DIFFUSION ───────────────────
\begin{frame}{GAN vs.\ Diffusion Models}
\begin{columns}[T]
\begin{column}{0.46\textwidth}
\begin{center}
\begin{tikzpicture}
\begin{axis}[width=6cm,height=5.5cm,
  symbolic x coords={Szybkość,Różnorodność,FID,Stabilność,Kontrola},
  xtick=data,xticklabel style={font=\tiny,rotate=25,anchor=east},
  ymin=0,ymax=10,ybar,bar width=0.32cm,
  legend style={font=\tiny,at={(0.97,0.97)},anchor=north east},
  title={\small Porównanie},title style={font=\small\bfseries},
  grid=major,grid style={gray!15}]
  \addplot[fill=ganaccent!80,draw=none]
    coordinates{(Szybkość,9)(Różnorodność,6)(FID,7)(Stabilność,5)(Kontrola,7)};
  \addlegendentry{GAN}
  \addplot[fill=ganteal!80,draw=none]
    coordinates{(Szybkość,3)(Różnorodność,9)(FID,9)(Stabilność,9)(Kontrola,8)};
  \addlegendentry{Diffusion}
\end{axis}
\end{tikzpicture}
\end{center}
\end{column}
\begin{column}{0.51\textwidth}
\vspace{0.4em}
\begin{tabular}{@{}lcc@{}}
\toprule & \textbf{GAN} & \textbf{Diffusion} \\ \midrule
Generowanie & 1 forward pass & 50–1000 kroków \\
Czas & $\sim$ms & sekundy \\
FID & dobry & lepszy \\
Mode collapse & tak & nie \\
Trening & trudny & stabilny \\
\bottomrule
\end{tabular}
\vspace{0.6em}
\abox{\textbf{Nisza GANów:} wszędzie gdzie kluczowa jest \textbf{szybkość} — super-resolution w grach (AMD FSR), real-time video, hybrydy GAN+Diffusion}
\end{column}
\end{columns}
\end{frame}

%% ─────────────────────── ETYKA ──────────────────────────────
\begin{frame}{Etyka i deepfake'i}
\begin{columns}[T]
\begin{column}{0.50\textwidth}
\textbf{Zagrożenia:}
\begin{itemize}
  \item Deepfake'i — szantaż, manipulacja polityczna
  \item Dezinformacja wizualna
  \item Bias modelu = bias w danych
  \item NCII (non-consensual intimate imagery)
\end{itemize}
\vspace{0.4em}
\textbf{Odpowiedzi:}
\begin{itemize}
  \item Detektory deepfake'ów
  \item Digital watermarking (C2PA)
\end{itemize}
\vspace{0.4em}
\abox{\textbf{EU AI Act (2024):} obowiązkowe oznaczanie treści generowanych przez AI}
\end{column}
\begin{column}{0.47\textwidth}
\begin{center}
\begin{tikzpicture}
\pie[color={ganaccent!80,ganblue!70,ganteal!60,gangray!50},
    text=legend,radius=2.1,style={font=\small}]{
  35/Dezinformacja,
  30/Deepfake wideo,
  20/Bias \& fairness,
  15/Prywatność}
\end{tikzpicture}
\end{center}
\end{column}
\end{columns}
\end{frame}

%% ─────────────────────── DEMO ───────────────────────────────
\begin{frame}{Demo — GAN na MNIST (PyTorch)}
\begin{columns}[T]
\begin{column}{0.50\textwidth}
\begin{block}{Architektura MLP-GAN}
\textbf{Generator:}
\[\mathbf{z}(100)\xrightarrow{\text{FC+BN+ReLU}}256\to512\to784\xrightarrow{\tanh}\hat{\mathbf{x}}\]
\textbf{Discriminator:}
\[\mathbf{x}(784)\xrightarrow{\text{FC+LReLU}}512\to256\to1\xrightarrow{\sigma}D(\mathbf{x})\]
\end{block}
\vspace{0.4em}
\abox{\textbf{Na co zwrócić uwagę:}\\
\texttt{.detach()} przy kroku D\\
Etykieta \texttt{1} przy kroku G\\
Normalizacja do $[-1,1]$ (Tanh)}
\end{column}
\begin{column}{0.47\textwidth}
\imgph[.95\linewidth]{Epoch 1 — szum wizualny}
\vspace{0.4em}
\imgph[.95\linewidth]{Epoch 50 — rozpoznawalne cyfry}
\end{column}
\end{columns}
\end{frame}

%% ─────────────────────── PODSUMOWANIE ───────────────────────
\begin{frame}{Podsumowanie}
\begin{columns}[T]
\begin{column}{0.55\textwidth}
\begin{tikzpicture}[
  it/.style={draw=ganblue!30,fill=ganlightbg,rounded corners=4pt,
    minimum width=7.2cm,minimum height=0.70cm,
    font=\small,inner sep=6pt,align=left},
  nm/.style={circle,fill=ganaccent,text=white,font=\bfseries\small,
    minimum size=0.50cm,inner sep=0}]
  \node[it](i1) at(0,0){Generator vs.\ Discriminator = czysta rywalizacja};
  \node[nm] at([xshift=6pt]i1.west){1};
  \node[it](i2) at(0,-0.88cm){$\min_G\max_D V(D,G)$ — gra dwuosobowa};
  \node[nm] at([xshift=6pt]i2.west){2};
  \node[it](i3) at(0,-1.76cm){WGAN: Wasserstein zamiast JS $\to$ stabilność};
  \node[nm] at([xshift=6pt]i3.west){3};
  \node[it](i4) at(0,-2.64cm){StyleGAN: rozplątana przestrzeń $\mathcal{W}$};
  \node[nm] at([xshift=6pt]i4.west){4};
  \node[it](i5) at(0,-3.52cm){CycleGAN: brak sparowanych danych};
  \node[nm] at([xshift=6pt]i5.west){5};
  \node[it](i6) at(0,-4.40cm){Nisza GANów: szybkość (ms) vs.\ Diffusion (s)};
  \node[nm] at([xshift=6pt]i6.west){6};
\end{tikzpicture}
\end{column}
\begin{column}{0.42\textwidth}
\begin{center}
\begin{tikzpicture}
\begin{axis}[width=5.8cm,height=5.6cm,
  title={\small Mapa wariantów GAN},title style={font=\small\bfseries},
  xlabel={\small kontrola outputu},ylabel={\small jakość obrazu},
  xmin=0,xmax=10,ymin=0,ymax=10,
  grid=major,grid style={gray!15},axis lines=left]
  \addplot[only marks,mark=*,mark size=5pt,ganblue]
    coordinates{(2,3)} node[right,font=\tiny]{GAN'14};
  \addplot[only marks,mark=*,mark size=5pt,ganaccent]
    coordinates{(3,6)} node[right,font=\tiny]{DCGAN};
  \addplot[only marks,mark=*,mark size=5pt,ganteal]
    coordinates{(7,5)} node[right,font=\tiny]{CGAN};
  \addplot[only marks,mark=*,mark size=5pt,ganblue!60]
    coordinates{(4,7)} node[right,font=\tiny]{WGAN};
  \addplot[only marks,mark=*,mark size=5pt,ganaccent!80]
    coordinates{(8,9)} node[right,font=\tiny]{StyleGAN};
  \addplot[only marks,mark=*,mark size=5pt,ganteal!70]
    coordinates{(6,6)} node[right,font=\tiny]{CycleGAN};
\end{axis}
\end{tikzpicture}
\end{center}
\end{column}
\end{columns}
\end{frame}

%% ─────────────────────── KONIEC ─────────────────────────────
{\setbeamertemplate{footline}{}
\begin{frame}[plain]
\vspace{2.5cm}
\begin{center}
  {\huge\bfseries\color{ganblue}Dziękujemy!}\\[0.8em]
  {\large Pytania?}\\[1.8em]
  \begin{tikzpicture}
    \node[draw=ganaccent,fill=ganaccent!10,rounded corners=8pt,
          inner sep=12pt,font=\large]{
      Maciej Leśniak \quad\textbullet\quad Bartłomiej Obrochta};
  \end{tikzpicture}\\[1.2em]
  {\small\color{gangray}Sztuczne Sieci Neuronowe $\cdot$ AGH Kraków $\cdot$ 2026}
\end{center}
\end{frame}}

\end{document}